\documentclass[11pt]{article}
\usepackage[T1]{fontenc}
\usepackage{textcomp}
\usepackage[margin=0.75in]{geometry}
\usepackage{amsmath,amssymb,amsthm}
\usepackage{array,booktabs}
\usepackage{hyperref}
\setlength{\parindent}{0pt}
\newtheorem{theorem}{Theorem}
\theoremstyle{definition}
\newtheorem{definition}{Definition}
\title{Flow Proof Artifact: prop-32-exterior-angle-sum}
\author{Generated by \texttt{flow doc proof}}
\date{Flow Proof Book}
\begin{document}
\maketitle
In a triangle, an exterior angle equals the sum of the two remote interior angles.\\[0.5em]
\textbf{Source.} euclid — Elements, Book I, Proposition 32\\[0.5em]
\bigskip
\noindent{\large \textbf{Derived fact 1.} \textit{In a triangle, an exterior angle equals the sum of the two remote interior angles} \hfill the Euclidean plane \cdot Euclid Book I \cdot Proposition 32: an exterior angle equals the sum of remote interior angles}\par
\smallskip\noindent\textit{Coordinate: the Euclidean plane \cdot Euclid Book I \cdot Proposition 32: an exterior angle equals the sum of remote interior angles}\par
\smallskip\noindent\textit{Source: euclid — Elements, Book I, Proposition 32}\par
\smallskip\noindent\textit{Built on: proposition 29: alternate angles are equal when a transversal crosses parallels, for Book I of the Elements in on the Euclidean plane, proposition 31: interior angles on one side sum to two right angles, for Book I of the Elements in on the Euclidean plane}\par
\medskip
\noindent\textbf{Goal.} In a triangle, an exterior angle equals the sum of the two remote interior angles\par
\noindent$\displaystyle \angle ACD = \angle CAB + \angle ABC$\par
\medskip
\noindent
\begin{tabular}{@{} >{\raggedright\arraybackslash}p{0.06\textwidth} >{\raggedright\arraybackslash}p{0.47\textwidth} >{\raggedleft\arraybackslash}p{0.41\textwidth} @{}}
\textbf{\#} & \textbf{Proof} & \textbf{Mathematics} \\
\midrule
\textbf{1.} & We prove that proposition 32: an exterior angle equals the sum of remote interior angles for Book I of the Elements in the Euclidean plane. & \\[0.45em]
\textbf{2.} & Let D = point\_on\_BC\_produced\_beyond\_C. & \\[0.45em]
\textbf{3.} & Let line\_through\_C\_parallel\_to\_AB = parallel\_at\_C\_to\_AB. & \\[0.45em]
\textbf{4.} & We invoke the derived fact: triangle\_ABC\_with\_exterior\_angle\_ACD. & \\[0.45em]
\textbf{5.} & We invoke the derived fact governing Book I of the Elements in the Euclidean plane: proposition 29: alternate angles are equal when a transversal crosses parallels, for Book I of the Elements in on the Euclidean plane. & \\[0.45em]
\textbf{6.} & We invoke the derived fact governing Book I of the Elements in the Euclidean plane: proposition 31: interior angles on one side sum to two right angles, for Book I of the Elements in on the Euclidean plane. & \\[0.45em]
\textbf{7.} & From step 2, step 3, step 4, step 5, and step 6, by the chain of results established in Book I (step 2, step 3, step 4, step 5, and step 6), we obtain angle ACD equals angle CAB plus angle ABC. Hence proven. & $\displaystyle \angle ACD = \angle CAB + \angle ABC$ \\[0.45em]
\bottomrule
\end{tabular}
\label{thm:Geometry-Euclid Book I-proposition_32_an_exterior_angle_equals_the_sum_of_remote_interior_angles}
\par\medskip\hrule
\end{document}
